\chapter{Concussion Tests}
\section*{What Is This Test?}
Concussion testing evaluates symptoms and neurologic function after a blow, bump, or jolt to the head or body. Concussion is a clinical diagnosis; no single routine scan or blood test confirms every concussion.
\section*{Alternative Names}
\begin{itemize}\item Mild Traumatic Brain Injury Evaluation\item mTBI Assessment\item Post-Concussion Assessment\end{itemize}
\section*{Principle}
Assessment combines injury history, symptom review, neurologic examination, balance and coordination testing, and cognitive assessment. CT is used when bleeding or another serious injury is suspected; MRI is reserved for selected cases.
\section*{Why This Test Is Ordered}
It is used after head trauma with headache, dizziness, confusion, memory difficulty, nausea, light sensitivity, slowed thinking, or balance problems. Worsening headache, repeated vomiting, seizure, weakness, unequal pupils, increasing confusion, or difficulty waking requires emergency evaluation.
\section*{Primary vs Secondary Test}
Clinical assessment is primary. Imaging is selected according to age, mechanism, examination, anticoagulant use, symptoms, and validated decision rules; routine imaging is not required for every uncomplicated concussion.
\section*{How This Test Is Performed}
The clinician documents the event, symptoms, orientation, memory, eye movements, strength, sensation, gait, balance, and coordination. Selected patients undergo CT or MRI.
\section*{What Preparation Is Needed}
No preparation is needed. Provide the injury history, medicines, anticoagulants, alcohol or drug exposure, loss of consciousness, and amnesia.
\section*{Understanding Results}
Normal imaging does not exclude concussion. Symptoms may evolve over hours or days and recovery is followed clinically with a gradual return to activities.
\section*{Follow-up Tests}
Follow-up may include repeat neurologic examination, vestibular or balance assessment, neuropsychological testing, vision evaluation, cervical-spine assessment, or specialist referral.
\section*{References}
\begin{enumerate}\item MedlinePlus. Concussion Tests.\item Centers for Disease Control and Prevention. Mild Traumatic Brain Injury and Concussion.\end{enumerate}
\clearpage\section*{Understanding Results and Follow-up}
The laboratory report should identify the specimen, method, units, reference interval or decision limit, and whether the result is positive, negative, abnormal, or inconclusive. Reference intervals vary with age, sex, pregnancy, specimen type, assay, and laboratory; a result outside a range is not automatically a diagnosis, and a result within a range does not always exclude disease. Discuss unexpected results with the ordering clinician, who can decide whether repeat or confirmatory testing, treatment, or referral is appropriate. Do not change medicines or delay urgent care based on an isolated result.

\section*{Regulatory Status and Authorization Date}This chapter describes a test category, not a specific commercial product. FDA, CDSCO, EU IVDR, and MHRA approval or registration dates therefore cannot be assigned without the assay manufacturer, product name, instrument, and intended use. For laboratory-developed tests, an individual agency approval date may not exist. See the Preface for the interpretation of regulatory status.
\clearpage
